{\color{ink}\bfseries 空时阻塞读、满时阻塞写}\par\vspace{1bp}以Linux阻塞式管道为例：空管道在仍有写端打开时等待；全部写端关闭且已有数据读尽后，read返回0，表示EOF。存在读端且一次写入的字节数0<n\(\le \)PIPE\_BUF时，要求整体原子写入，剩余空间不足n字节仍会等待，因此“有空间即可写”还须核对本次请求大小。fork后若保留不用的写端，读者可能一直等不到EOF。
